% Источники: 07_Программные_требования/DOC-06.2_Сессии_задач_и_планировщик.md, 07_Программные_требования/Контракты_сессий.md
% Генератор: docs/build_pdf.py
% Шаблон печатного представления. Содержимое подставляет docs/build_pdf.py.
\documentclass[12pt,a4paper]{article}
\usepackage{iftex}
\ifPDFTeX
  \usepackage[utf8]{inputenc}
  \usepackage[T2A]{fontenc}
  \usepackage{newtxtext,newtxmath}
\else
  \usepackage{fontspec}
  \IfFontExistsTF{Times New Roman}
    {\setmainfont{Times New Roman}}
    {\setmainfont{Libertinus Serif}}
  \IfFontExistsTF{Arial}{\setsansfont{Arial}}{\setsansfont{Latin Modern Sans}}
  \IfFontExistsTF{Menlo}{\setmonofont{Menlo}[Scale=MatchLowercase]}
    {\setmonofont{Latin Modern Mono}[Scale=MatchLowercase]}
\fi
\usepackage[russian]{babel}
\usepackage{geometry}
\geometry{left=25mm,right=15mm,top=20mm,bottom=20mm}
\usepackage{setspace}
\onehalfspacing
\usepackage{indentfirst}
\setlength{\parindent}{1.25cm}
\setlength{\parskip}{2pt}
\setlength{\emergencystretch}{3em}
\usepackage{graphicx,booktabs,longtable,array,tabularx,calc,etoolbox}
\usepackage{enumitem}
\setlist[itemize]{leftmargin=1.25cm,itemsep=2pt,topsep=4pt}
\setlist[enumerate]{leftmargin=1.25cm,itemsep=2pt,topsep=4pt}
\newcommand{\tightlist}{\setlength{\itemsep}{0pt}\setlength{\parskip}{0pt}}
\AtBeginEnvironment{longtable}{\small\singlespacing}
\setlength{\LTpre}{8pt}
\setlength{\LTpost}{8pt}
\setlength{\tabcolsep}{5pt}
\setlength{\LTleft}{0pt}
\setlength{\LTright}{0pt}
\usepackage{xurl}
\usepackage[unicode,colorlinks=true,linkcolor=black,urlcolor=blue]{hyperref}
\usepackage{bookmark}
\hypersetup{pdftitle={Сессии задач и планировщик},pdfauthor={Анохин Г.А.},pdfsubject={Криптотренажёр mephi-12}}
\urlstyle{same}
\pagestyle{plain}
\setcounter{secnumdepth}{3}
\setcounter{tocdepth}{2}

\begin{document}
\hypersetup{pageanchor=false}
\begin{titlepage}
  \thispagestyle{empty}
  \centering
  {\large НИЯУ МИФИ\par}
  \vspace*{3.5cm}
  {\LARGE\bfseries Сессии задач и планировщик\par}
  \vspace{0.5cm}
  {\large DOC-06.2\par}
  \vspace{0.6cm}
  {\large Платформа для изучения и практики курса по криптографии\par}
  {\large Криптотренажёр mephi-12\par}
  \vspace{0.7cm}
  {\large Версия 1.0\par}
  {\large Редакция от 20.11.2026\par}
  \vspace{1cm}
  \begin{tabularx}{\textwidth}{@{}p{3.4cm}Xp{4cm}@{}}
    \toprule
    Роль & ФИО & Подпись / Дата \\
    \midrule
    Разработал & Анохин Г.А. & \\
    \addlinespace[6pt]
    Согласовано & & \\
    \addlinespace[6pt]
    Утверждено & & \\
    \bottomrule
  \end{tabularx}
  \vfill
  {\large Москва, 2026\par}
\end{titlepage}
\hypersetup{pageanchor=true}

\section*{Лист регистрации изменений}
\addcontentsline{toc}{section}{Лист регистрации изменений}
\begingroup\small\singlespacing
\noindent
\begin{tabularx}{\textwidth}{@{}p{1.2cm}p{2cm}p{2.5cm}Xp{1.8cm}@{}}
\toprule
Версия & Дата & Автор & Содержание изменения & Утвердил \\
\midrule
1.0 & 20.11.2026 & Анохин Г.А. & Переработаны требования к жизненному циклу сессии, восстановлению и контрактам операций. & \\
\bottomrule
\end{tabularx}
\endgroup
\clearpage
\begingroup\small\singlespacing
\tableofcontents
\endgroup
\clearpage

\hypertarget{doc-062-сессии-задач-и-планировщик}{}

{\def\LTcaptype{none} % do not increment counter
\begin{longtable}[]{@{}
  >{\raggedright\arraybackslash}p{(\linewidth - 2\tabcolsep) * \real{0.3200}}
  >{\raggedright\arraybackslash}p{(\linewidth - 2\tabcolsep) * \real{0.6800}}@{}}
\toprule\noalign{}
\begin{minipage}[b]{\linewidth}\raggedright
Реквизит
\end{minipage} & \begin{minipage}[b]{\linewidth}\raggedright
Значение
\end{minipage} \\
\midrule\noalign{}
\endhead
\bottomrule\noalign{}
\endlastfoot
Редакция & 1.0 от 20.11.2026 \\
Ответственный & Анохин Г.А. \\
Основания & DOC-05, системные требования к сессиям и интеграции. DOC-06, общие программные требования \\
Приложение & Контракты сессий задач \\
\end{longtable}
}

\section{Назначение и состав}\label{1-ux43dux430ux437ux43dux430ux447ux435ux43dux438ux435-ux438-ux441ux43eux441ux442ux430ux432}

\hypertarget{1-назначение-и-состав}{}

Личная сессия --- выданный одному студенту набор конкретных задач с общим сроком решения. Она возникает, когда студент приступает к опубликованному контрольному мероприятию своей группы или начинает самостоятельную практику. Каждая задача получает собственное условие. Предметный обработчик проверяет ответ студента по этому условию и возвращает вердикт «верно» или «неверно». Подсистема сессий решает, можно ли принять ответ, сохраняет вердикт и подводит итог набора. Планировщик --- фоновая часть программы, которая по серверным часам находит просроченные задачи и закрывает их, даже если студент отключился.

Правила этой подсистемы обязательны. В разделах 2--4 каждому правилу присвоен ID. Системные основания и контрольные случаи приведены в {\hyperlink{6-трассировка-и-описательные-проверки}{таблице трассировки}}. Формы запросов и ответов определяет {\hyperlink{doc-062-контракты-сессий-задач}{приложение}}. Смежные обязанности определены в {\href{DOC-06\_Программные\_требования.pdf\#doc-06-сводные-программные-требования}{общих программных требованиях}}.

\subsection{Данные и их жизненный цикл}\label{ux434ux430ux43dux43dux44bux435-ux438-ux438ux445-ux436ux438ux437ux43dux435ux43dux43dux44bux439-ux446ux438ux43aux43b}

\hypertarget{данные-и-их-жизненный-цикл}{}

UUID --- уникальный идентификатор, по которому сервер находит конкретную запись. У пользователя, мероприятия, сессии, каждой задачи и каждой операции свои идентификаторы. Слово «операция» ниже означает обработку одного поступившего запроса с решением. Наименование операции и смысл её параметров приведены в {\hyperlink{операции}{контракте сессий}}.

{\def\LTcaptype{none} % do not increment counter
\begin{longtable}[]{@{}
  >{\raggedright\arraybackslash}p{(\linewidth - 2\tabcolsep) * \real{0.3200}}
  >{\raggedright\arraybackslash}p{(\linewidth - 2\tabcolsep) * \real{0.6800}}@{}}
\toprule\noalign{}
\begin{minipage}[b]{\linewidth}\raggedright
Объект
\end{minipage} & \begin{minipage}[b]{\linewidth}\raggedright
Что хранится и для чего используется
\end{minipage} \\
\midrule\noalign{}
\endhead
\bottomrule\noalign{}
\endlastfoot
Сессия & UUID, владелец, \texttt{assessmentId} контрольного мероприятия либо \texttt{null} для практики, серверное время создания, срок подачи решений \texttt{deadlineAt}, состав и состояние. Все задачи одного набора имеют общий срок \\
Задача & UUID, связь с сессией, выданное условие, тип, состояние и зафиксированный предметный результат, если он получен \\
Принятый ответ & Входные данные и срок конкретной операции. По этой записи обработчик может продолжить прерванную проверку, а планировщик определяет, какие задачи пока нельзя закрывать \\
\end{longtable}
}

Для задачи \texttt{NEW} означает «итог ещё не определён», \texttt{SOLVED} --- «зафиксировано верное решение», \texttt{FAILED} --- «зафиксировано неверное решение или истёк срок». Во время проверки задача остаётся \texttt{NEW}. Для всей сессии \texttt{NEW} означает наличие хотя бы одной незавершённой задачи. Только после завершения последней задачи определяется {\hyperlink{swr-ses-018-общая-функция-итога-сессии}{итог сессии}}. Отметка о принятом ответе относится к операции. Завершённая предметная попытка возникает при записи вердикта. Прекращение операции из-за ошибки формата или технического сбоя само по себе попытку не расходует.

\subsection{Время решения и время обработки}\label{ux432ux440ux435ux43cux44f-ux440ux435ux448ux435ux43dux438ux44f-ux438-ux432ux440ux435ux43cux44f-ux43eux431ux440ux430ux431ux43eux442ux43aux438}

\hypertarget{время-решения-и-время-обработки}{}

Серверное время создания сессии обозначено \texttt{t0}. Для самостоятельной практики срок подачи \texttt{D\ =\ t0\ +\ 20\ минут}. Для контроля \texttt{D\ =\ min(t0\ +\ 20\ минут,\ closesAt)}, где \texttt{closesAt} --- момент закрытия опубликованной работы. В ответе сервера этот общий для задач срок передаётся как \texttt{deadlineAt}. Момент получения сервером полного запроса обозначен \texttt{t\_in}. Запрос своевременен, если \texttt{t\_in\ \textless{}\ D}. Для него действует отдельный предел фиксации результата \texttt{L =\allowbreak{} t\_\allowbreak{}in + 60 секунд}. В момент \texttt{L} записывать вердикт уже поздно. Плановый период фонового контроля тоже равен 60 секундам. Это частота запуска планировщика, она не изменяет ни срок подачи, ни предел обработки. {\hyperlink{согласование-параллельных-событий}{Пример на часах}} показывает случай, когда оба срока наступают в разное время.

\section{Создание сессии}\label{2-ux441ux43eux437ux434ux430ux43dux438ux435-ux441ux435ux441ux441ux438ux438}

\hypertarget{2-создание-сессии}{}

Самостоятельная практика использует шаблон набора задач --- упорядоченный перечень типов, доступный группе. Для контроля состав задан опубликованным преподавателем мероприятием. Проверку доступности и выбор источника состава описывают требования ниже. Одна транзакция хранилища объединяет запись сессии и всех её задач. Успешная фиксация делает их доступными вместе, а откат отменяет изменения этого создания.

\subsection{SWR-SES-001. Получение пользователя}\label{swr-ses-001-ux43fux43eux43bux443ux447ux435ux43dux438ux435-ux43fux43eux43bux44cux437ux43eux432ux430ux442ux435ux43bux44f}

\hypertarget{swr-ses-001-получение-пользователя}{}

\hypertarget{swr-ses-001}{}

При создании набора и при отправке решения обработчик должен получить личность пользователя от компонента проверки доступа до обращения к данным сессии. Если доступ не подтверждён, обработка запроса прекращается. Идентификатор в данных клиента не предоставляет права действовать от имени другого владельца.

\subsection{SWR-SES-002. Выбор доступного шаблона}\label{swr-ses-002-ux432ux44bux431ux43eux440-ux434ux43eux441ux442ux443ux43fux43dux43eux433ux43e-ux448ux430ux431ux43bux43eux43dux430}

\hypertarget{swr-ses-002-выбор-доступного-шаблона}{}

\hypertarget{swr-ses-002}{}

Для самостоятельной практики компонент создания должен установить группу по подтверждённому университетскому профилю студента и найти среди шаблонов этой группы включённый шаблон с запрошенным именем. Из него компонент берёт состав набора. Для пользователя без группы, при отсутствии имени в группе и при отключённом шаблоне результат один --- \texttt{CONFIGURATION\_\allowbreak{}NOT\_\allowbreak{}FOUND}. Сессия при этом не создаётся.

Для контроля клиент передаёт \texttt{assessmentId}. Компонент находит опубликованное мероприятие подтверждённой группы пользователя. Чужое, неопубликованное и отсутствующее мероприятие дают \texttt{404 ASSESSMENT\_\allowbreak{}NOT\_\allowbreak{}FOUND}. До открытия (\texttt{opensAt}) старт отклоняется с \texttt{409 ASSESSMENT\_\allowbreak{}NOT\_\allowbreak{}OPEN}, с момента закрытия (\texttt{closesAt}) --- с \texttt{409 ASSESSMENT\_\allowbreak{}CLOSED}. Внутри окна сервер проверяет активную сессию по {\hyperlink{swr-ses-003}{SWR-SES-003}}. После завершения предыдущего прохождения этой работы новый набор выдаётся только при \texttt{repeatAllowed\ =\ true}. Если повторы запрещены, сервер возвращает \texttt{409 ASSESSMENT\_\allowbreak{}RETAKE\_\allowbreak{}DISABLED} с UUID прежней сессии в \texttt{details.\allowbreak{}sessionId}. При разрешённом старте формируется новый набор по закреплённому составу с тем же \texttt{assessmentId}. Ранее выданные наборы сохраняются. У практики \texttt{assessmentId\ =\ null}.

\subsection{SWR-SES-003. Защита от двух активных сессий}\label{swr-ses-003-ux437ux430ux449ux438ux442ux430-ux43eux442-ux434ux432ux443ux445-ux430ux43aux442ux438ux432ux43dux44bux445-ux441ux435ux441ux441ux438ux439}

\hypertarget{swr-ses-003-защита-от-двух-активных-сессий}{}

\hypertarget{swr-ses-003}{}

Компонент создания должен в одной транзакции заблокировать строку пользователя, учесть просрочку прежнего набора и проверить, осталась ли у этого пользователя сессия \texttt{NEW}. Внутри открытого окна контроля и при создании практики сессия \texttt{NEW} исключает выдачу второго набора независимо от \texttt{repeatAllowed}. Решение о повторном старте и запись нового прохождения выполняются под той же блокировкой. Второй параллельный запрос после снятия блокировки получает \texttt{ACTIVE\_\allowbreak{}SESSION\_\allowbreak{}EXISTS} с идентификатором действующей сессии.

Даже после общего срока прежняя сессия остаётся активной, пока у неё есть задача со своевременно полученным ответом в обработке. Порядок завершения такой задачи определяют {\hyperlink{swr-ses-015}{SWR-SES-015}} и {\hyperlink{swr-ses-020}{SWR-SES-020}}.

\subsection{SWR-SES-004. Формирование задач}\label{swr-ses-004-ux444ux43eux440ux43cux438ux440ux43eux432ux430ux43dux438ux435-ux437ux430ux434ux430ux447}

\hypertarget{swr-ses-004-формирование-задач}{}

\hypertarget{swr-ses-004}{}

Для каждой позиции выбранного состава компонент создания должен запросить условие у модуля указанного типа. В сессии задачи следуют в порядке, заданном шаблоном. При повторе типа каждой позиции соответствует отдельная задача со своим UUID и выданным условием. Новая сессия и все её задачи получают состояние \texttt{NEW}. Последующие правки шаблона не изменяют состав набора и условия выданных задач. Если генератор для заданного типа отсутствует, создание завершается техническим отказом.

Модуль может сформировать условие вычислением или выбрать подготовленный вариант. Для позиции типа \texttt{CUSTOM\_\allowbreak{}TEXT} он случайно выбирает одну из заданных преподавателем пар «условие --- эталон». Выдача закрепляет выбор за задачей личной сессии. Студенту возвращается только условие. Остальные варианты и эталон выбранного варианта в ответ не включаются. Повторное чтение возвращает то же условие.

\subsection{SWR-SES-005. Сохранение созданного набора}\label{swr-ses-005-ux441ux43eux445ux440ux430ux43dux435ux43dux438ux435-ux441ux43eux437ux434ux430ux43dux43dux43eux433ux43e-ux43dux430ux431ux43eux440ux430}

\hypertarget{swr-ses-005-сохранение-созданного-набора}{}

\hypertarget{swr-ses-005}{}

Сессию и все её задачи компонент создания должен сохранить одной транзакцией PostgreSQL. Сообщить об успехе можно после фиксации полного состава. При ошибке генератора или записи транзакция откатывается. У пользователя не появляется частичный набор, который занимал бы место активной сессии.

\subsection{SWR-SES-006. Сохранение времени создания}\label{swr-ses-006-ux441ux43eux445ux440ux430ux43dux435ux43dux438ux435-ux432ux440ux435ux43cux435ux43dux438-ux441ux43eux437ux434ux430ux43dux438ux44f}

\hypertarget{swr-ses-006-сохранение-времени-создания}{}

\hypertarget{swr-ses-006}{}

При создании подсистема должна сохранить серверное время \texttt{t0} и общий срок \texttt{deadlineAt}. Для практики это \texttt{t0\ +\ 20\ минут}. Для контроля сервер выбирает более ранний из двух моментов --- \texttt{t0\ +\ 20\ минут} или \texttt{closesAt}. Чтение сессии, повторный запрос, ошибка и перезапуск не меняют исходных временных границ. Время создания отдельных задач на них не влияет.

\section{Обработка ответа}\label{3-ux43eux431ux440ux430ux431ux43eux442ux43aux430-ux43eux442ux432ux435ux442ux430}

\hypertarget{3-обработка-ответа}{}

Проверки идут в два этапа. Допуск устанавливает, вправе ли этот пользователь решать задачу сейчас. Проверка формата устанавливает, можно ли прочитать решение по схеме выданного задания. Только затем предметный обработчик оценивает сам ответ. При конкуренции запросов блокировка даёт одному обработчику возможность работать с общей записью, пока другие ждут её освобождения. Это защищает запись единственного вердикта.

\subsection{SWR-SES-007. Регистрация времени поступления}\label{swr-ses-007-ux440ux435ux433ux438ux441ux442ux440ux430ux446ux438ux44f-ux432ux440ux435ux43cux435ux43dux438-ux43fux43eux441ux442ux443ux43fux43bux435ux43dux438ux44f}

\hypertarget{swr-ses-007-регистрация-времени-поступления}{}

\hypertarget{swr-ses-007}{}

Приёмный обработчик должен зафиксировать \texttt{t\_in} в момент получения полного запроса. До очереди он сверяет подтверждённого отправителя с владельцем задачи. Для отсутствующей или чужой задачи возвращается \texttt{TASK\_\allowbreak{}NOT\_\allowbreak{}FOUND}, операция для этой задачи не открывается. Ответ владельца сохраняется с исходной отметкой по {\hyperlink{swr-ses-021}{SWR-SES-021}}. Регистрация своевременного запроса владельца согласуется с планировщиком по \texttt{t\_in}, чтобы планировщик учитывал его до закрытия задачи. После перезапуска обработчик продолжает проверки допуска и формата с прежним \texttt{t\_in}. Время на устройстве клиента и начало математического вычисления на допуск не влияют.

\subsection{SWR-SES-008. Проверка допуска}\label{swr-ses-008-ux43fux440ux43eux432ux435ux440ux43aux430-ux434ux43eux43fux443ux441ux43aux430}

\hypertarget{swr-ses-008-проверка-допуска}{}

\hypertarget{swr-ses-008}{}

До вызова проверяющего подсистема должна установить по сохранённым данным, что задача принадлежит сессии текущего пользователя, задача и сессия ещё допускают ответ, а запрос получен до \texttt{D}. Для отсутствующей и чужой задачи возвращается \texttt{TASK\_\allowbreak{}NOT\_\allowbreak{}FOUND}. После установления владельца сначала решается вопрос о сроке. При \texttt{t\_in\ \textgreater{}=\ D} возвращается \texttt{SESSION\_\allowbreak{}EXPIRED}, даже если задача уже имеет итог или планировщик ещё не обработал запись.

У своевременного запроса к завершённой задаче или сессии результат --- \texttt{TASK\_\allowbreak{}FINALIZED}. При незавершённой проверке другого допустимого ответа к той же задаче действует {\hyperlink{swr-ses-012}{SWR-SES-012}}.

\subsection{SWR-SES-009. Проверка формата}\label{swr-ses-009-ux43fux440ux43eux432ux435ux440ux43aux430-ux444ux43eux440ux43cux430ux442ux430}

\hypertarget{swr-ses-009-проверка-формата}{}

\hypertarget{swr-ses-009}{}

До предметного вычисления обработчик должен разобрать JSON и проверить обязательные поля и их типы по схеме сохранённой задачи. При ошибке формы он возвращает \texttt{INVALID\_\allowbreak{}REQUEST}. При формально допустимой записи вычислений проверяющий получает ответ, даже если расчёт окажется неверным. Предел размера тела устанавливает {\hyperlink{swr-ses-022}{SWR-SES-022}}.

\subsection{SWR-SES-010. Передача профильной проверке}\label{swr-ses-010-ux43fux435ux440ux435ux434ux430ux447ux430-ux43fux440ux43eux444ux438ux43bux44cux43dux43eux439-ux43fux440ux43eux432ux435ux440ux43aux435}

\hypertarget{swr-ses-010-передача-профильной-проверке}{}

\hypertarget{swr-ses-010}{}

Подсистема должна выбрать проверяющий по типу сохранённой задачи и передать ему исходное выданное условие вместе с принятым ответом. Поля пользователя не изменяют условие или выбор проверяющего. Неизвестный тип и техническая недоступность обработчика завершают операцию без предметного вердикта.

Для \texttt{BACKPACK}, \texttt{EL\_GAMAL}, \texttt{CUSTOM\_\allowbreak{}TEXT} и подключаемых типов действуют одни правила времени и учёта попытки. Модуль \texttt{CUSTOM\_\allowbreak{}TEXT} извлекает эталон закреплённого варианта. После удаления пробелов по краям обеих строк он сравнивает их с учётом регистра и всех внутренних пробелов. Равенство даёт положительный вердикт. Схему тела задаёт {\hyperlink{пользовательское-текстовое-задание}{приложение}}.

\subsection{SWR-SES-011. Сохранение вердикта}\label{swr-ses-011-ux441ux43eux445ux440ux430ux43dux435ux43dux438ux435-ux432ux435ux440ux434ux438ux43aux442ux430}

\hypertarget{swr-ses-011-сохранение-вердикта}{}

\hypertarget{swr-ses-011}{}

При получении вердикта подсистема должна установить, что операция ещё открыта и время её фиксации не истекло. Дальше она записывает предметный результат с UUID и типом задачи, устанавливает задаче \texttt{SOLVED} для верного решения либо \texttt{FAILED} для неверного, пересчитывает состояние сессии и закрывает операцию. Все изменения фиксируются вместе. Подтверждение клиенту следует за фиксацией.

\subsection{SWR-SES-012. Согласование одновременных обращений}\label{swr-ses-012-ux441ux43eux433ux43bux430ux441ux43eux432ux430ux43dux438ux435-ux43eux434ux43dux43eux432ux440ux435ux43cux435ux43dux43dux44bux445-ux43eux431ux440ux430ux449ux435ux43dux438ux439}

\hypertarget{swr-ses-012-согласование-одновременных-обращений}{}

\hypertarget{swr-ses-012}{}

Подсистема должна атомарно предоставить право на предметную проверку одному допустимому ответу на задачу. Два зарегистрированных запроса с разными UUID не дают права на два одновременных вычисления по одной задаче. При ещё открытой проверке второй своевременный допустимый запрос получает \texttt{CHECK\_\allowbreak{}IN\_\allowbreak{}PROGRESS}, при уже зафиксированном результате до общего срока --- \texttt{TASK\_\allowbreak{}FINALIZED}. В обоих случаях тело ранее принятого ответа сохраняется.

Обработчик вычисляет вердикт вне транзакции записи. Для фиксации он открывает короткую транзакцию и блокирует по порядку строку сессии, строку задачи и запись операции. После получения блокировок он проверяет актуальные данные. При параллельном завершении разных задач одной сессии итог рассчитывается по их согласованным состояниям под блокировкой сессии.

\subsection{SWR-SES-013. Сохранение итогового состояния}\label{swr-ses-013-ux441ux43eux445ux440ux430ux43dux435ux43dux438ux435-ux438ux442ux43eux433ux43eux432ux43eux433ux43e-ux441ux43eux441ux442ux43eux44fux43dux438ux44f}

\hypertarget{swr-ses-013-сохранение-итогового-состояния}{}

\hypertarget{swr-ses-013}{}

Подсистема должна отклонить новую отправку решения для задачи \texttt{SOLVED} или \texttt{FAILED} до вызова проверяющего. При повторе запроса сохранённый результат и состояние задачи остаются прежними. Владелец получает результат отдельной операцией чтения.

\subsection{SWR-SES-014. Отказ без расхода попытки}\label{swr-ses-014-ux43eux442ux43aux430ux437-ux431ux435ux437-ux440ux430ux441ux445ux43eux434ux430-ux43fux43eux43fux44bux442ux43aux438}

\hypertarget{swr-ses-014-отказ-без-расхода-попытки}{}

\hypertarget{swr-ses-014}{}

При ошибке формата и при техническом отказе без зафиксированного вердикта обработчик должен вернуть причину отказа и сохранить предметную попытку. До \texttt{D} задача остаётся \texttt{NEW}, если у неё нет другого основания для завершения. Пользователь может подать исправленный ответ с новым \texttt{t\_in}. Когда срок наступил, к оставшейся без своевременной обработки задаче применяется {\hyperlink{swr-ses-017}{SWR-SES-017}}.

Если запись вердикта откатилась и обработчик сообщает клиенту \texttt{500 INTERNAL\_\allowbreak{}ERROR}, он должен сначала сохранить отмену операции. После доставки отказа прежний результат уже не может быть записан. При обрыве связи или остановке процесса судьбу ответа определяет хранилище. Открытая запись подлежит восстановлению по {\hyperlink{swr-ses-021}{SWR-SES-021}}, а зафиксированный результат доступен для чтения.

\subsection{SWR-SES-015. Завершение своевременной проверки}\label{swr-ses-015-ux437ux430ux432ux435ux440ux448ux435ux43dux438ux435-ux441ux432ux43eux435ux432ux440ux435ux43cux435ux43dux43dux43eux439-ux43fux440ux43eux432ux435ux440ux43aux438}

\hypertarget{swr-ses-015-завершение-своевременной-проверки}{}

\hypertarget{swr-ses-015}{}

Для запроса владельца с \texttt{t\_in\ \textless{}\ D} обработчик и планировщик должны сохранить возможность пройти допуск и проверку формата, а после успешного допуска --- записать предметный вердикт, в том числе после \texttt{D}. Это правило действует и при закрытии контрольного мероприятия. Вердикт должен быть зафиксирован строго до \texttt{L}. Защита от просрочки действует с регистрации своевременного ответа подтверждённого владельца и включает ожидание в пределах {\hyperlink{swr-ses-020}{SWR-SES-020}}. Обращение другого пользователя с UUID задачи не препятствует её закрытию. Если запрос отклонён или проверка завершилась техническим отказом, связанная с ним защита снимается. При этом учитываются остальные своевременные запросы владельца, если они ещё обрабатываются. После \texttt{D} новый запрос уже не допускается.

\section{Контроль срока и итог}\label{4-ux43aux43eux43dux442ux440ux43eux43bux44c-ux441ux440ux43eux43aux430-ux438-ux438ux442ux43eux433}

\hypertarget{4-контроль-срока-и-итог}{}

Контроль срока завершает задачу без предметного вердикта, когда прошёл {\hyperlink{время-решения-и-время-обработки}{общий срок подачи ответа}} и нет своевременной обработки. После каждого завершения пересчитывается состояние всего набора. Истечение времени и неверное решение имеют одно состояние задачи \texttt{FAILED}, но различаются по {\hyperlink{чтение-после-потери-ответа}{сохранённому результату}}.

\subsection{SWR-SES-016. Фоновый контроль}\label{swr-ses-016-ux444ux43eux43dux43eux432ux44bux439-ux43aux43eux43dux442ux440ux43eux43bux44c}

\hypertarget{swr-ses-016-фоновый-контроль}{}

\hypertarget{swr-ses-016}{}

Планировщик должен обрабатывать просроченные задачи с плановым периодом 60 секунд. Такой период задан требованием ТЗ {\href{../02\_Техническое\_задание/Техническое\_задание.pdf\#ts-req-func-session-time}{\texttt{TS.\allowbreak{}REQ.\allowbreak{}FUNC.\allowbreak{}SESSION.\allowbreak{}TIME}}}. Независимо от последнего запуска планировщика приёмный обработчик проверяет каждый новый ответ относительно \texttt{D}.

\subsection{SWR-SES-017. Завершение по просрочке}\label{swr-ses-017-ux437ux430ux432ux435ux440ux448ux435ux43dux438ux435-ux43fux43e-ux43fux440ux43eux441ux440ux43eux447ux43aux435}

\hypertarget{swr-ses-017-завершение-по-просрочке}{}

\hypertarget{swr-ses-017}{}

По достижении \texttt{D} подсистема должна закрывать по просрочке только задачу \texttt{NEW}, для которой нет своевременного запроса в обработке. Её состояние меняется на \texttt{FAILED}, поле \texttt{result} остаётся равным \texttt{null}. После изменения пересчитывается итог набора. Планировщик сохраняет прежние вердикты и учитывает защиту незавершённых операций по {\hyperlink{swr-ses-015}{SWR-SES-015}}.

\subsection{SWR-SES-018. Общая функция итога сессии}\label{swr-ses-018-ux43eux431ux449ux430ux44f-ux444ux443ux43dux43aux446ux438ux44f-ux438ux442ux43eux433ux430-ux441ux435ux441ux441ux438ux438}

\hypertarget{swr-ses-018-общая-функция-итога-сессии}{}

\hypertarget{swr-ses-018}{}

После предметной проверки и после закрытия задачи по сроку подсистема должна применять одинаковое правило к состояниям всех задач сессии.

{\def\LTcaptype{none} % do not increment counter
\begin{longtable}[]{@{}
  >{\raggedright\arraybackslash}p{(\linewidth - 2\tabcolsep) * \real{0.3200}}
  >{\raggedright\arraybackslash}p{(\linewidth - 2\tabcolsep) * \real{0.6800}}@{}}
\toprule\noalign{}
\begin{minipage}[b]{\linewidth}\raggedright
Состояния задач
\end{minipage} & \begin{minipage}[b]{\linewidth}\raggedright
Состояние сессии
\end{minipage} \\
\midrule\noalign{}
\endhead
\bottomrule\noalign{}
\endlastfoot
Хотя бы одна задача остаётся \texttt{NEW} & \texttt{NEW}, включая сочетание \texttt{NEW} и \texttt{FAILED} \\
Все задачи завершены и каждая имеет \texttt{SOLVED} & \texttt{SOLVED} \\
Все задачи завершены, среди них есть \texttt{FAILED} & \texttt{FAILED} \\
\end{longtable}
}

\subsection{SWR-SES-019. Сохранность данных}\label{swr-ses-019-ux441ux43eux445ux440ux430ux43dux43dux43eux441ux442ux44c-ux434ux430ux43dux43dux44bux445}

\hypertarget{swr-ses-019-сохранность-данных}{}

\hypertarget{swr-ses-019}{}

Между запусками приложения хранилище должно сохранять UUID и принадлежность сессии и задач, состав, условия и типы задач, \texttt{t0}, установленный срок подачи \texttt{deadlineAt}, текущие состояния и зафиксированные результаты. После запуска приложение использует те же условия и сроки. Для ответа, по которому ещё нет вердикта, используется запись операции по {\hyperlink{swr-ses-021}{SWR-SES-021}}.

\subsection{SWR-SES-020. Предел обработки}\label{swr-ses-020-ux43fux440ux435ux434ux435ux43b-ux43eux431ux440ux430ux431ux43eux442ux43aux438}

\hypertarget{swr-ses-020-предел-обработки}{}

\hypertarget{swr-ses-020}{}

Вердикт можно записать только при времени фиксации строго меньше \texttt{t\_\allowbreak{}in + 60 секунд}. По достижении этой границы подсистема должна закрыть незавершённую операцию с \texttt{CHECK\_\allowbreak{}TIMEOUT}, отбросить поздний вердикт и сохранить неиспользованную попытку. В один интервал входят ожидание очереди, вычисление, фиксация и восстановление. Истечение интервала обработки не меняет \texttt{D}. Состояние задачи после отказа определяется {\hyperlink{swr-ses-014}{SWR-SES-014}} и {\hyperlink{swr-ses-017}{SWR-SES-017}}.

\subsection{SWR-SES-021. Запись принятой операции}\label{swr-ses-021-ux437ux430ux43fux438ux441ux44c-ux43fux440ux438ux43dux44fux442ux43eux439-ux43eux43fux435ux440ux430ux446ux438ux438}

\hypertarget{swr-ses-021-запись-принятой-операции}{}

\hypertarget{swr-ses-021}{}

После сверки владельца по {\hyperlink{swr-ses-007}{SWR-SES-007}} и до очереди подсистема должна записать в устойчивое хранилище UUID операции, подтверждённого пользователя, UUID собственной задачи из запроса, полный ответ, \texttt{t\_in}, предел \texttt{t\_\allowbreak{}in + 60 секунд} и признак незавершённости. Запись запроса нужна для учёта времени и восстановления. Право начать предметную проверку возникает после {\hyperlink{swr-ses-008}{SWR-SES-008}} и {\hyperlink{swr-ses-009}{SWR-SES-009}} и предоставляется на задачу по {\hyperlink{swr-ses-012}{SWR-SES-012}}. Блокировка при получении записи не позволяет двум обработчикам одновременно продолжить одну операцию.

После перезапуска приложение должно найти открытые операции, в том числе ещё ожидавшие в очереди. До исходного предела обработчик продолжает работу по сохранённым данным с учётом правил допуска и формата. При истёкшем пределе действует {\hyperlink{swr-ses-020}{SWR-SES-020}}. Закрытие операции, запись вердикта и новые состояния задачи и сессии принадлежат одной транзакции. Повторный запуск после фиксации использует записанный результат.

\subsection{SWR-SES-022. Форматы и категории API}\label{swr-ses-022-ux444ux43eux440ux43cux430ux442ux44b-ux438-ux43aux430ux442ux435ux433ux43eux440ux438ux438-api}

\hypertarget{swr-ses-022-форматы-и-категории-api}{}

\hypertarget{swr-ses-022}{}

Интерфейс должен применять список обязательных полей, точные типы, представление времени и целых чисел, ограничение размера и соответствие HTTP-статусов кодам из {\hyperlink{doc-062-контракты-сессий-задач}{приложения}}. Создание принимает ровно один параметр из \texttt{configuration} и \texttt{assessmentId}. Отсутствие обоих, их сочетание, повтор или ошибка значения дают \texttt{400 INVALID\_\allowbreak{}REQUEST}. Ошибки параметров, структуры тела и его размера прекращают обработку до предметной проверки. Неподдерживаемый тип задания в сохранённом шаблоне даёт технический отказ с откатом создания.

\subsection{SWR-SES-023. Чтение сессии и результата}\label{swr-ses-023-ux447ux442ux435ux43dux438ux435-ux441ux435ux441ux441ux438ux438-ux438-ux440ux435ux437ux443ux43bux44cux442ux430ux442ux430}

\hypertarget{swr-ses-023-чтение-сессии-и-результата}{}

\hypertarget{swr-ses-023}{}

По \texttt{GET /\allowbreak{}api/\allowbreak{}tasks/\allowbreak{}session/\allowbreak{}\{sessionId\}} подсистема должна предоставить владельцу его сессию и задачи с сохранёнными результатами, в том числе после закрытия мероприятия и после начала нового разрешённого прохождения. Чужая и отсутствующая сессии дают \texttt{SESSION\_\allowbreak{}NOT\_\allowbreak{}FOUND}. Перед выдачей подсистема приводит состояния к правилам просрочки и учитывает своевременные обработки. Чтение не вызывает предметного вычисления и не расходует попытку.

\section{Временные границы}\label{5-ux432ux440ux435ux43cux435ux43dux43dux44bux435-ux433ux440ux430ux43dux438ux446ux44b}

\hypertarget{5-временные-границы}{}

\subsection{Согласование параллельных событий}\label{ux441ux43eux433ux43bux430ux441ux43eux432ux430ux43dux438ux435-ux43fux430ux440ux430ux43bux43bux435ux43bux44cux43dux44bux445-ux441ux43eux431ux44bux442ux438ux439}

\hypertarget{согласование-параллельных-событий}{}

Пусть сервер создал сессию практики в 10:00:00 и получил полный ответ в 10:19:59. В 10:20:00 планировщик видит ещё открытую задачу. Регистрация запроса защищает её от просрочки, даже если проверяющий пока не приступил к вычислению. До 10:20:59 обработчик может зафиксировать вердикт. Если к 10:20:59 фиксации нет, операция завершается с \texttt{CHECK\_\allowbreak{}TIMEOUT}, а задача --- по просрочке без предметного результата. В той же сессии другой полный запрос, впервые полученный ровно в 10:20:00, уже опоздал.

Контрольное мероприятие закрывается в 11:00:00. Студент начинает в 10:57:00, поэтому его срок \texttt{D} --- 11:00:00. Для ответа, полностью полученного в 10:59:59, запись вердикта разрешена строго до 11:00:59. В 11:00:00 новый ответ владельца получает \texttt{410 SESSION\_\allowbreak{}EXPIRED}, а запрос на новый старт по мероприятию --- \texttt{409 ASSESSMENT\_\allowbreak{}CLOSED}. Прежнюю сессию можно читать в обоих случаях.

\subsection{Переходы состояния задачи}\label{ux43fux435ux440ux435ux445ux43eux434ux44b-ux441ux43eux441ux442ux43eux44fux43dux438ux44f-ux437ux430ux434ux430ux447ux438}

\hypertarget{переходы-состояния-задачи}{}

{\def\LTcaptype{none} % do not increment counter
\begin{longtable}[]{@{}
  >{\raggedright\arraybackslash}p{(\linewidth - 2\tabcolsep) * \real{0.3200}}
  >{\raggedright\arraybackslash}p{(\linewidth - 2\tabcolsep) * \real{0.6800}}@{}}
\toprule\noalign{}
\begin{minipage}[b]{\linewidth}\raggedright
Событие для задачи в состоянии \texttt{NEW}
\end{minipage} & \begin{minipage}[b]{\linewidth}\raggedright
Состояние и результат, доступные владельцу
\end{minipage} \\
\midrule\noalign{}
\endhead
\bottomrule\noalign{}
\endlastfoot
Зафиксирован положительный вердикт & \texttt{SOLVED}, результат с \texttt{isOk\ =\ true} \\
Зафиксирован отрицательный вердикт & \texttt{FAILED}, результат с \texttt{isOk\ =\ false} и пояснением \\
До общего срока произошла ошибка формата или технический отказ без вердикта & \texttt{NEW}, \texttt{result\ =\ null}, если нет другого основания закрытия \\
Общий срок наступил во время своевременной обработки & \texttt{NEW}, пока обработка может дать вердикт \\
Общий срок наступил, своевременных обработок нет & \texttt{FAILED}, \texttt{result\ =\ null} \\
\end{longtable}
}

\subsection{Проверка длительной обработки}\label{ux43fux440ux43eux432ux435ux440ux43aux430-ux434ux43bux438ux442ux435ux43bux44cux43dux43eux439-ux43eux431ux440ux430ux431ux43eux442ux43aux438}

\hypertarget{проверка-длительной-обработки}{}

При разборе временного случая фиксируют три момента по серверным часам --- получение полного запроса, общий срок сессии и запись вердикта. Начало вычисления само по себе ещё не определяет исход. Например, проверяющий может вернуть верный ответ за секунду до предела, а транзакция записи получит блокировку уже на границе. Такой результат подпадает под {\hyperlink{swr-ses-020}{SWR-SES-020}}. Клиент сможет прочитать состояние задачи по её сессии.

\section{Трассировка и описательные проверки}\label{6-ux442ux440ux430ux441ux441ux438ux440ux43eux432ux43aux430-ux438-ux43eux43fux438ux441ux430ux442ux435ux43bux44cux43dux44bux435-ux43fux440ux43eux432ux435ux440ux43aux438}

\hypertarget{6-трассировка-и-описательные-проверки}{}

Системные источники приведены в {\href{../04\_Системные\_требования/DOC-05\_Системные\_требования.pdf\#doc-05-требования-к-сессиям-задач}{требованиях к сессиям}} и {\href{../04\_Системные\_требования/DOC-05\_Системные\_требования.pdf\#doc-05-общесистемная-интеграция}{правилах интеграции}}. В последнем столбце указано, какое наблюдение позволяет проверить обязанность.

{\def\LTcaptype{none} % do not increment counter
\begin{longtable}[]{@{}
  >{\raggedright\arraybackslash}p{(\linewidth - 4\tabcolsep) * \real{0.2131}}
  >{\raggedright\arraybackslash}p{(\linewidth - 4\tabcolsep) * \real{0.3063}}
  >{\raggedright\arraybackslash}p{(\linewidth - 4\tabcolsep) * \real{0.4807}}@{}}
\toprule\noalign{}
\begin{minipage}[b]{\linewidth}\raggedright
Требование
\end{minipage} & \begin{minipage}[b]{\linewidth}\raggedright
Системные источники
\end{minipage} & \begin{minipage}[b]{\linewidth}\raggedright
Контрольный случай и ожидаемый признак
\end{minipage} \\
\midrule\noalign{}
\endhead
\bottomrule\noalign{}
\endlastfoot
{\hyperlink{swr-ses-001}{SWR-SES-001}} & {\href{../04\_Системные\_требования/DOC-05\_Системные\_требования.pdf\#sr-ses-003}{SR-SES-003}}, {\href{../04\_Системные\_требования/DOC-05\_Системные\_требования.pdf\#sr-sys-001}{SR-SYS-001}} & Передан чужой ID. Операция доступна только от имени подтверждённого пользователя \\
{\hyperlink{swr-ses-002}{SWR-SES-002}} & {\href{../04\_Системные\_требования/DOC-05\_Системные\_требования.pdf\#sr-ses-001}{SR-SES-001}}, {\href{../04\_Системные\_требования/DOC-05\_Системные\_требования.pdf\#sr-sys-003}{SR-SYS-003}} & Нет группы или включённого шаблона. Пользователь не получил набор \\
{\hyperlink{swr-ses-003}{SWR-SES-003}} & {\href{../04\_Системные\_требования/DOC-05\_Системные\_требования.pdf\#sr-ses-002}{SR-SES-002}} & Два одновременных запроса создания одного пользователя дают единственную активную сессию \\
{\hyperlink{swr-ses-004}{SWR-SES-004}} & {\href{../04\_Системные\_требования/DOC-05\_Системные\_требования.pdf\#sr-ses-001}{SR-SES-001}}, {\href{../04\_Системные\_требования/DOC-05\_Системные\_требования.pdf\#sr-sys-002}{SR-SYS-002}}, {\href{../04\_Системные\_требования/DOC-05\_Системные\_требования.pdf\#sr-sys-003}{SR-SYS-003}}, {\href{../04\_Системные\_требования/DOC-05\_Системные\_требования.pdf\#sr-sys-004}{SR-SYS-004}} & В шаблоне повторяется тип. В наборе две разные задачи этого типа на своих позициях. Изменение шаблона после выдачи не меняет набора \\
{\hyperlink{swr-ses-005}{SWR-SES-005}} & {\href{../04\_Системные\_требования/DOC-05\_Системные\_требования.pdf\#sr-ses-001}{SR-SES-001}}, {\href{../04\_Системные\_требования/DOC-05\_Системные\_требования.pdf\#sr-sys-005}{SR-SYS-005}} & Ошибка при формировании последней задачи. Первая задача и незавершённая сессия не видны пользователю \\
{\hyperlink{swr-ses-006}{SWR-SES-006}} & {\href{../04\_Системные\_требования/DOC-05\_Системные\_требования.pdf\#sr-ses-006}{SR-SES-006}}, {\href{../04\_Системные\_требования/DOC-05\_Системные\_требования.pdf\#sr-sys-006}{SR-SYS-006}}, {\href{../04\_Системные\_требования/DOC-05\_Системные\_требования.pdf\#sr-sys-007}{SR-SYS-007}} & После повторного чтения и перезапуска срок практики по-прежнему отсчитывается от создания. Срок контрольной сессии также ограничивает \texttt{closesAt} опубликованной работы \\
{\hyperlink{swr-ses-007}{SWR-SES-007}} & {\href{../04\_Системные\_требования/DOC-05\_Системные\_требования.pdf\#sr-ses-011}{SR-SES-011}}, {\href{../04\_Системные\_требования/DOC-05\_Системные\_требования.pdf\#sr-sys-006}{SR-SYS-006}} & Перезапуск застал запрос в очереди. Сохранились сам ответ и время его полного получения \\
{\hyperlink{swr-ses-008}{SWR-SES-008}} & {\href{../04\_Системные\_требования/DOC-05\_Системные\_требования.pdf\#sr-ses-003}{SR-SES-003}}, {\href{../04\_Системные\_требования/DOC-05\_Системные\_требования.pdf\#sr-ses-008}{SR-SES-008}}, {\href{../04\_Системные\_требования/DOC-05\_Системные\_требования.pdf\#sr-ses-011}{SR-SES-011}} & Для чужой задачи и позднего запроса к собственной задаче получены соответствующие отказы до вычисления \\
{\hyperlink{swr-ses-009}{SWR-SES-009}} & {\href{../04\_Системные\_требования/DOC-05\_Системные\_требования.pdf\#sr-ses-009}{SR-SES-009}}, {\href{../04\_Системные\_требования/DOC-05\_Системные\_требования.pdf\#sr-sys-009}{SR-SYS-009}} & Отсутствует обязательное поле. Сервер сообщает об ошибке формата без вердикта \\
{\hyperlink{swr-ses-010}{SWR-SES-010}} & {\href{../04\_Системные\_требования/DOC-05\_Системные\_требования.pdf\#sr-ses-004}{SR-SES-004}}, {\href{../04\_Системные\_требования/DOC-05\_Системные\_требования.pdf\#sr-sys-004}{SR-SYS-004}} & На проверку поступили условие и тип именно той задачи, UUID которой указан в запросе \\
{\hyperlink{swr-ses-011}{SWR-SES-011}} & {\href{../04\_Системные\_требования/DOC-05\_Системные\_требования.pdf\#sr-ses-005}{SR-SES-005}}, {\href{../04\_Системные\_требования/DOC-05\_Системные\_требования.pdf\#sr-ses-007}{SR-SES-007}}, {\href{../04\_Системные\_требования/DOC-05\_Системные\_требования.pdf\#sr-sys-005}{SR-SYS-005}} & После подтверждения видны записанный вердикт, соответствующее состояние задачи и пересчитанный итог набора \\
{\hyperlink{swr-ses-012}{SWR-SES-012}} & {\href{../04\_Системные\_требования/DOC-05\_Системные\_требования.pdf\#sr-ses-008}{SR-SES-008}} & Два ответа к одной задаче перекрылись во времени. Зафиксирован один предметный результат \\
{\hyperlink{swr-ses-013}{SWR-SES-013}} & {\href{../04\_Системные\_требования/DOC-05\_Системные\_требования.pdf\#sr-ses-008}{SR-SES-008}}, {\href{../04\_Системные\_требования/DOC-05\_Системные\_требования.pdf\#sr-sys-005}{SR-SYS-005}} & Повтор после фиксации получил отказ. Состояние и первый результат сохранились \\
{\hyperlink{swr-ses-014}{SWR-SES-014}} & {\href{../04\_Системные\_требования/DOC-05\_Системные\_требования.pdf\#sr-ses-009}{SR-SES-009}}, {\href{../04\_Системные\_требования/DOC-05\_Системные\_требования.pdf\#sr-ses-010}{SR-SES-010}}, {\href{../04\_Системные\_требования/DOC-05\_Системные\_требования.pdf\#sr-sys-009}{SR-SYS-009}} & После технического отказа до срока можно подать новый ответ, \texttt{D} сохранился \\
{\hyperlink{swr-ses-015}{SWR-SES-015}} & {\href{../04\_Системные\_требования/DOC-05\_Системные\_требования.pdf\#sr-ses-011}{SR-SES-011}}, {\href{../04\_Системные\_требования/DOC-05\_Системные\_требования.pdf\#sr-sys-006}{SR-SYS-006}} & После общего срока своевременный ответ по-прежнему может дать вердикт до своего предела \\
{\hyperlink{swr-ses-016}{SWR-SES-016}} & {\href{../04\_Системные\_требования/DOC-05\_Системные\_требования.pdf\#sr-ses-006}{SR-SES-006}}, {\href{../04\_Системные\_требования/DOC-05\_Системные\_требования.pdf\#sr-sys-006}{SR-SYS-006}} & Планировщик задержан. Новый ответ в момент общего срока отклонён \\
{\hyperlink{swr-ses-017}{SWR-SES-017}} & {\href{../04\_Системные\_требования/DOC-05\_Системные\_требования.pdf\#sr-ses-006}{SR-SES-006}}, {\href{../04\_Системные\_требования/DOC-05\_Системные\_требования.pdf\#sr-ses-011}{SR-SES-011}} & По сроку закрыты только задачи без итога и без своевременной обработки \\
{\hyperlink{swr-ses-018}{SWR-SES-018}} & {\href{../04\_Системные\_требования/DOC-05\_Системные\_требования.pdf\#sr-ses-007}{SR-SES-007}}, {\href{../04\_Системные\_требования/DOC-05\_Системные\_требования.pdf\#sr-sys-005}{SR-SYS-005}}, {\href{../04\_Системные\_требования/DOC-05\_Системные\_требования.pdf\#sr-sys-006}{SR-SYS-006}} & Одна задача \texttt{FAILED}, другая ещё \texttt{NEW}. Сессия получает \texttt{FAILED} только после завершения второй \\
{\hyperlink{swr-ses-019}{SWR-SES-019}} & {\href{../04\_Системные\_требования/DOC-05\_Системные\_требования.pdf\#sr-sys-002}{SR-SYS-002}}, {\href{../04\_Системные\_требования/DOC-05\_Системные\_требования.pdf\#sr-sys-007}{SR-SYS-007}} & После перезапуска совпадают владелец, состав, времена, условия и сохранённые вердикты \\
{\hyperlink{swr-ses-020}{SWR-SES-020}} & {\href{../04\_Системные\_требования/DOC-05\_Системные\_требования.pdf\#sr-ses-010}{SR-SES-010}}, {\href{../04\_Системные\_требования/DOC-05\_Системные\_требования.pdf\#sr-ses-012}{SR-SES-012}} & При попытке фиксации ровно на границе обработки вердикт отброшен \\
{\hyperlink{swr-ses-021}{SWR-SES-021}} & {\href{../04\_Системные\_требования/DOC-05\_Системные\_требования.pdf\#sr-ses-008}{SR-SES-008}}, {\href{../04\_Системные\_требования/DOC-05\_Системные\_требования.pdf\#sr-ses-013}{SR-SES-013}}, {\href{../04\_Системные\_требования/DOC-05\_Системные\_требования.pdf\#sr-sys-005}{SR-SYS-005}}, {\href{../04\_Системные\_требования/DOC-05\_Системные\_требования.pdf\#sr-sys-007}{SR-SYS-007}} & При восстановлении доступны прежние входные данные. До предела допускается одна фиксация, после предела --- технический отказ \\
{\hyperlink{swr-ses-022}{SWR-SES-022}} & {\href{../04\_Системные\_требования/DOC-05\_Системные\_требования.pdf\#sr-ses-009}{SR-SES-009}}, {\href{../04\_Системные\_требования/DOC-05\_Системные\_требования.pdf\#sr-sys-008}{SR-SYS-008}}, {\href{../04\_Системные\_требования/DOC-05\_Системные\_требования.pdf\#sr-sys-009}{SR-SYS-009}} & Повтор параметра, превышение размера и неверная схема различаются в соответствии с таблицей ошибок \\
{\hyperlink{swr-ses-023}{SWR-SES-023}} & {\href{../04\_Системные\_требования/DOC-05\_Системные\_требования.pdf\#sr-ses-014}{SR-SES-014}}, {\href{../04\_Системные\_требования/DOC-05\_Системные\_требования.pdf\#sr-sys-002}{SR-SYS-002}}, {\href{../04\_Системные\_требования/DOC-05\_Системные\_требования.pdf\#sr-sys-007}{SR-SYS-007}}, {\href{../04\_Системные\_требования/DOC-05\_Системные\_требования.pdf\#sr-sys-008}{SR-SYS-008}} & После потери ответа владелец читает сохранённый результат. Чужой пользователь получает тот же отказ, что и для отсутствующей сессии \\
\end{longtable}
}

\section{История редакций}\label{ux438ux441ux442ux43eux440ux438ux44f-ux440ux435ux434ux430ux43aux446ux438ux439}

\hypertarget{история-редакций}{}

{\def\LTcaptype{none} % do not increment counter
\begin{longtable}[]{@{}
  >{\raggedright\arraybackslash}p{(\linewidth - 4\tabcolsep) * \real{0.2623}}
  >{\raggedright\arraybackslash}p{(\linewidth - 4\tabcolsep) * \real{0.2623}}
  >{\raggedright\arraybackslash}p{(\linewidth - 4\tabcolsep) * \real{0.4755}}@{}}
\toprule\noalign{}
\begin{minipage}[b]{\linewidth}\raggedright
Дата
\end{minipage} & \begin{minipage}[b]{\linewidth}\raggedright
Версия
\end{minipage} & \begin{minipage}[b]{\linewidth}\raggedright
Изменение
\end{minipage} \\
\midrule\noalign{}
\endhead
\bottomrule\noalign{}
\endlastfoot
05.11.2026 & 0.8 & Описаны создание сессии, обработка ответов и планировщик \\
07.11.2026 & 0.9 & Определены контракты, предел проверки и восстановление операций \\
20.11.2026 & 1.0 & Переработаны требования и приложение. Уточнено изложение порядка фиксации, временных границ и чтения результата \\
\end{longtable}
}

\section{DOC-06.2. Контракты сессий задач}\label{doc-062-ux43aux43eux43dux442ux440ux430ux43aux442ux44b-ux441ux435ux441ux441ux438ux439-ux437ux430ux434ux430ux447}

\hypertarget{doc-062-контракты-сессий-задач}{}

Приложение к {\hyperlink{doc-062-сессии-задач-и-планировщик}{требованиям подсистемы сессий и планировщика}}. Редакция 1.0 от 20.11.2026. Ответственный --- Анохин Г.А.

В этом приложении «клиент» означает программу, через которую студент работает с платформой. Клиент начинает личную сессию, отправляет решение одной задачи и читает состояние набора. Контракт описывает запрос и возможные ответы сервера для каждого действия. Представление сессии, задачи или результата для обмена называют объектом передачи данных (DTO). Контракт определяет поля этого объекта и допустимые значения. Полное внутреннее состояние приложения остаётся на сервере.

\subsection{Операции}\label{ux43eux43fux435ux440ux430ux446ux438ux438}

\hypertarget{операции}{}

В записи запроса первое слово --- метод HTTP, дальше указан путь. Метод \texttt{POST} используется для создания сессии, \texttt{PUT} --- для отправки решения, \texttt{GET} --- для чтения. Заполнители в фигурных скобках заменяются конкретными значениями. UUID --- уникальный идентификатор объекта. Код 201 сообщает о создании ресурса, 200 --- об успешной обработке запроса. Конкретную причину отказа определяют по {\hyperlink{коды-и-представление-ошибок}{таблице ошибок}}.

Все три операции требуют подтверждённой личности и роли \texttt{USER}, дающей студенту доступ к собственной работе. Сервер определяет владельца по {\href{DOC-06\_Программные\_требования.pdf\#swr-auth-003-проверка-токена-и-роли}{проверенному токену доступа}}. Данные передаются в JSON --- текстовом формате объектов с именованными полями и массивов значений. Для JSON используются кодировка UTF-8 и тип содержимого \texttt{application/\allowbreak{}json}.

{\def\LTcaptype{none} % do not increment counter
\begin{longtable}[]{@{}
  >{\raggedright\arraybackslash}p{(\linewidth - 4\tabcolsep) * \real{0.2320}}
  >{\raggedright\arraybackslash}p{(\linewidth - 4\tabcolsep) * \real{0.4182}}
  >{\raggedright\arraybackslash}p{(\linewidth - 4\tabcolsep) * \real{0.3498}}@{}}
\toprule\noalign{}
\begin{minipage}[b]{\linewidth}\raggedright
Запрос
\end{minipage} & \begin{minipage}[b]{\linewidth}\raggedright
Данные от клиента
\end{minipage} & \begin{minipage}[b]{\linewidth}\raggedright
Успешный ответ
\end{minipage} \\
\midrule\noalign{}
\endhead
\bottomrule\noalign{}
\endlastfoot
\texttt{POST /\allowbreak{}api/\allowbreak{}tasks/\allowbreak{}session} & Ровно один параметр. \texttt{assessmentId} --- UUID мероприятия из записи в расписании студента, \texttt{configuration} --- имя шаблона из списка для его группы & \texttt{201\ Created}. Созданная сессия со всеми задачами \\
\texttt{PUT /\allowbreak{}api/\allowbreak{}tasks/\allowbreak{}submit} & Один параметр запроса \texttt{statementId} с UUID задачи и тело с решением & \texttt{200\ OK}. Сохранённый вердикт для указанной задачи \\
\texttt{GET /\allowbreak{}api/\allowbreak{}tasks/\allowbreak{}session/\allowbreak{}\{sessionId\}} & UUID сессии в пути & \texttt{200\ OK}. Собственная сессия с состояниями и результатами задач \\
\end{longtable}
}

По \texttt{assessmentId} сервер находит опубликованное мероприятие группы пользователя и проверяет серверное время. Новая сессия допустима с момента открытия \texttt{opensAt} включительно до момента закрытия \texttt{closesAt} исключительно. Внутри окна при действующей сессии \texttt{NEW} новая не создаётся. Завершённое прохождение можно повторить, только если преподаватель выбрал \texttt{repeatAllowed\ =\ true} до публикации. Повтор получает собственные UUID сессии и задач. Прежнюю сессию владелец может прочитать по её UUID и после закрытия. Успешная выдача закрепляет за сессией идентификатор этого мероприятия. По имени шаблона создаётся практика без связи с контролем. Одновременная передача \texttt{configuration} и \texttt{assessmentId}, отсутствие обоих параметров, повтор или неверная запись значения дают \texttt{400 INVALID\_\allowbreak{}REQUEST}.

Имя шаблона в параметре \texttt{configuration} сравнивается с учётом регистра. Сервер ищет шаблон среди включённых для группы пользователя. Для отправки решения отсутствие, повтор или неверная запись \texttt{statementId} также дают \texttt{400 INVALID\_\allowbreak{}REQUEST}.

На тело решения отводится 1 МиБ (1 048 576 байт). Размер определяется до разбора JSON. При большем размере возвращается \texttt{413 PAYLOAD\_\allowbreak{}TOO\_\allowbreak{}LARGE}. Пустое тело, ошибка JSON, отсутствие обязательного поля, недопустимый тип или \texttt{null} в обязательном поле дают \texttt{400 INVALID\_\allowbreak{}REQUEST}. В обоих случаях предметная попытка сохраняется.

\subsection{Представление сессии и задачи}\label{ux43fux440ux435ux434ux441ux442ux430ux432ux43bux435ux43dux438ux435-ux441ux435ux441ux441ux438ux438-ux438-ux437ux430ux434ux430ux447ux438}

\hypertarget{представление-сессии-и-задачи}{}

Сессия --- личный набор студента, а задача --- один конкретный экземпляр задания внутри него. Первый объект ответа описывает сессию. В его массиве \texttt{problems} находятся объекты задач, по одному на каждую выданную позицию. Значение \texttt{null} обозначает отсутствие значения, например ещё не записанный результат. Подробнее о {\hyperlink{данные-и-их-жизненный-цикл}{состояниях задачи и сессии}}.

{\def\LTcaptype{none} % do not increment counter
\begin{longtable}[]{@{}
  >{\raggedright\arraybackslash}p{(\linewidth - 2\tabcolsep) * \real{0.3200}}
  >{\raggedright\arraybackslash}p{(\linewidth - 2\tabcolsep) * \real{0.6800}}@{}}
\toprule\noalign{}
\begin{minipage}[b]{\linewidth}\raggedright
Поле сессии
\end{minipage} & \begin{minipage}[b]{\linewidth}\raggedright
Тип и значение
\end{minipage} \\
\midrule\noalign{}
\endhead
\bottomrule\noalign{}
\endlastfoot
\texttt{id} & UUID созданной сессии \\
\texttt{assessmentId} & UUID опубликованного мероприятия для контрольной сессии, \texttt{null} для самостоятельной практики \\
\texttt{createdAt} & Серверное время создания набора. Дата и время по всемирному координированному времени (UTC) в формате RFC 3339 с миллисекундами и суффиксом \texttt{Z}, например \texttt{2026-\allowbreak{}11-\allowbreak{}20T12:00:00.\allowbreak{}000Z} \\
\texttt{deadlineAt} & Срок подачи ответа для всех задач, в том же формате UTC, что и \texttt{createdAt}. Для практики --- через 20 минут после создания, для контроля --- более ранний из \texttt{createdAt\ +\ 20\ минут} и \texttt{closesAt}. Новый ответ в этот момент уже опоздал. \\
\texttt{state} & Состояние набора. Принимает \texttt{NEW}, \texttt{SOLVED} или \texttt{FAILED} \\
\texttt{userId} & UUID владельца \\
\texttt{problems} & Массив задач именно этой сессии \\
\end{longtable}
}

{\def\LTcaptype{none} % do not increment counter
\begin{longtable}[]{@{}
  >{\raggedright\arraybackslash}p{(\linewidth - 2\tabcolsep) * \real{0.3200}}
  >{\raggedright\arraybackslash}p{(\linewidth - 2\tabcolsep) * \real{0.6800}}@{}}
\toprule\noalign{}
\begin{minipage}[b]{\linewidth}\raggedright
Поле задачи
\end{minipage} & \begin{minipage}[b]{\linewidth}\raggedright
Тип и значение
\end{minipage} \\
\midrule\noalign{}
\endhead
\bottomrule\noalign{}
\endlastfoot
\texttt{id} & UUID, который клиент передаёт в \texttt{statementId} \\
\texttt{statement} & Строка с JSON выданного условия \\
\texttt{type} & Сохранённый тип задачи. Определяет схему решения \\
\texttt{state} & Состояние задачи из того же перечня \texttt{NEW}, \texttt{SOLVED}, \texttt{FAILED} \\
\texttt{sessionId} & UUID сессии, к которой принадлежит задача \\
\texttt{result} & Объект предметного результата либо \texttt{null}, если вердикт не был зафиксирован \\
\end{longtable}
}

\subsubsection{Смысл связей в передаваемых данных}\label{ux441ux43cux44bux441ux43b-ux441ux432ux44fux437ux435ux439-ux432-ux43fux435ux440ux435ux434ux430ux432ux430ux435ux43cux44bux445-ux434ux430ux43dux43dux44bux445}

\hypertarget{смысл-связей-в-передаваемых-данных}{}

Для запроса чтения клиент использует \texttt{id} сессии. Для отправки решения он выбирает одну запись из \texttt{problems} и берёт \texttt{id} задачи. Ответ на решение содержит этот же UUID в поле \texttt{problemId}. Так клиент может сопоставить ответ со строкой задачи, даже если в наборе несколько задач одного типа.

Для всех задач набора применяется срок из \texttt{deadlineAt}. Срок определяется на сервере, часовой пояс клиента на него не влияет. Сервер определяет владельца задачи по её связи с сессией. Полученные из клиента идентификаторы не заменяют проверку владельца.

\subsection{Предметный ответ}\label{ux43fux440ux435ux434ux43cux435ux442ux43dux44bux439-ux43eux442ux432ux435ux442}

\hypertarget{предметный-ответ}{}

Каждый подключённый тип задаёт схему решения. Для \texttt{BACKPACK} и \texttt{EL\_GAMAL} передают расчётные величины, для \texttt{CUSTOM\_\allowbreak{}TEXT} --- текст ответа. При создании набор с неподдерживаемым типом отклоняется целиком с \texttt{500 INTERNAL\_\allowbreak{}ERROR}. При отправке решения схему диктует сохранённый тип задачи. Клиент не может передать новое условие или выбрать другого проверяющего полем своего ответа.

Для больших целых используется \textbf{десятичная строка}. Допустимы \texttt{"0"} и строка из цифр с ненулевой первой цифрой, перед которой для отрицательного целого стоит знак минуса. Записи \texttt{"-0"}, \texttt{"+2"}, \texttt{"02"}, \texttt{"2.0"}, \texttt{"2e3"} и пробелы не допускаются. Строковое представление сохраняет точное значение при обмене. Предметный проверяющий отдельно оценивает математические связи между корректно записанными значениями.

\subsubsection{Ранцевая задача}\label{ux440ux430ux43dux446ux435ux432ux430ux44f-ux437ux430ux434ux430ux447ux430}

\hypertarget{ранцевая-задача}{}

Студент передаёт параметры {\href{DOC-06.3\_Ранцевые\_криптосистемы.pdf\#как-устроена-ранцевая-задача}{ранцевого шифрования}}. Сообщение указывает, какие веса нужно сложить. Лёгкий и тяжёлый ранцы --- два связанных списка весов. Названия полей ниже описывают части одного решения.

В решении \texttt{BACKPACK} обязательны все поля таблицы. Ни одно из них не принимает \texttt{null}.

{\def\LTcaptype{none} % do not increment counter
\begin{longtable}[]{@{}
  >{\raggedright\arraybackslash}p{(\linewidth - 2\tabcolsep) * \real{0.3200}}
  >{\raggedright\arraybackslash}p{(\linewidth - 2\tabcolsep) * \real{0.6800}}@{}}
\toprule\noalign{}
\begin{minipage}[b]{\linewidth}\raggedright
Поле
\end{minipage} & \begin{minipage}[b]{\linewidth}\raggedright
Представление и смысл
\end{minipage} \\
\midrule\noalign{}
\endhead
\bottomrule\noalign{}
\endlastfoot
\texttt{power} & Положительное целое JSON-число, размерность задачи \\
\texttt{type} & Строка с вариантом ранца, совпадающим с вариантом выданного условия \\
\texttt{message} & Массив логических значений \texttt{true} и \texttt{false}, исходное сообщение \\
\texttt{lightBackpack} & Массив десятичных строк, элементы лёгкого ранца \\
\texttt{omega} & Десятичная строка, множитель преобразования \\
\texttt{hardBackpack} & Массив десятичных строк, элементы тяжёлого ранца \\
\texttt{encodedMessage} & Десятичная строка, зашифрованное сообщение \\
\texttt{decodedMessage} & Массив логических значений, восстановленное сообщение \\
\texttt{module} & Десятичная строка, модуль преобразования \\
\texttt{reverseOmega} & Десятичная строка, обратный множитель \\
\end{longtable}
}

\subsubsection{Эль-Гамаль}\label{ux44dux43bux44c-ux433ux430ux43cux430ux43bux44c}

\hypertarget{эль-гамаль}{}

В этом учебном шифре простое число \texttt{p} задаёт модуль, то есть число, по которому берут остатки, а \texttt{g} --- основание для возведения в степень. По заданному закрытому показателю \texttt{x} вычисляют открытый ключ \texttt{y}. Для шифрования сообщения \texttt{m} используется отдельный одноразовый показатель \texttt{k}. Результат шифрования состоит из двух чисел \texttt{c1} и \texttt{c2}. Проверяются и эти числа, и сообщение, полученное после расшифрования. Закрытый показатель выдан в составе учебного условия, чтобы студент мог выполнить весь расчёт.

В решении \texttt{EL\_GAMAL} пять обязательных полей. Каждое содержит десятичную строку. Значение \texttt{null} не допускается.

{\def\LTcaptype{none} % do not increment counter
\begin{longtable}[]{@{}
  >{\raggedright\arraybackslash}p{(\linewidth - 2\tabcolsep) * \real{0.3200}}
  >{\raggedright\arraybackslash}p{(\linewidth - 2\tabcolsep) * \real{0.6800}}@{}}
\toprule\noalign{}
\begin{minipage}[b]{\linewidth}\raggedright
Поле
\end{minipage} & \begin{minipage}[b]{\linewidth}\raggedright
Значение
\end{minipage} \\
\midrule\noalign{}
\endhead
\bottomrule\noalign{}
\endlastfoot
\texttt{y} & Открытый ключ \\
\texttt{k} & Одноразовый параметр шифрования \\
\texttt{c1} & Первая составляющая шифротекста \\
\texttt{c2} & Вторая составляющая шифротекста \\
\texttt{decryptedMessage} & Восстановленное сообщение \\
\end{longtable}
}

\subsubsection{Учебный пример решения Эль-Гамаля}\label{ux443ux447ux435ux431ux43dux44bux439-ux43fux440ux438ux43cux435ux440-ux440ux435ux448ux435ux43dux438ux44f-ux44dux43bux44c-ux433ux430ux43cux430ux43bux44f}

\hypertarget{учебный-пример-решения-эль-гамаля}{}

Для условия с \texttt{p\ =\ 23}, \texttt{g\ =\ 5}, \texttt{x\ =\ 6}, \texttt{m\ =\ 13} пользователь выбрал \texttt{k\ =\ 3}. Тело корректного ответа имеет следующий вид.

\begin{verbatim}
{
  "y": "8",
  "k": "3",
  "c1": "10",
  "c2": "9",
  "decryptedMessage": "13"
}
\end{verbatim}

Замена \texttt{"13"} числом \texttt{13} нарушит тип поля и даст ошибку формата. Строка \texttt{"12"} типу соответствует, но математически неверна для этого условия. В таком случае клиент получит предметный результат с \texttt{isOk\ =\ false}.

\subsubsection{Пользовательское текстовое задание}\label{ux43fux43eux43bux44cux437ux43eux432ux430ux442ux435ux43bux44cux441ux43aux43eux435-ux442ux435ux43aux441ux442ux43eux432ux43eux435-ux437ux430ux434ux430ux43dux438ux435}

\hypertarget{пользовательское-текстовое-задание}{}

В редакторе преподаватель вводит условие варианта как обычный текст и отдельно задаёт правильный ответ. У выданных задач всех видов одинаковые внешние поля. Поле \texttt{statement} содержит строку с сериализованным JSON, то есть с текстовой записью значения в формате JSON. У типа \texttt{CUSTOM\_\allowbreak{}TEXT} внутреннее значение само является строкой. В ответе сервера возможен такой фрагмент объекта задачи.

\begin{verbatim}
{"type":"CUSTOM_TEXT","statement":"\"Сколько бит в байте?\""}
\end{verbatim}

Клиент разбирает внешний объект, берёт строку из \texttt{statement} и разбирает её как JSON. Теперь он может показать студенту текст «Сколько бит в байте?» без кавычек. Ответ передаётся как обычный JSON-объект \texttt{\{"answer":"8"\}}; второй слой сериализации для него не нужен.

При публикации фиксируется редакция вариантов из задания преподавателя. При выдаче задачи \texttt{CUSTOM\_\allowbreak{}TEXT} модуль случайно выбирает одну пару из этой редакции для указанной позиции. Студент получает только условие выбранного варианта. Два студента могут получить одинаковый вариант. Прочие варианты и эталонный ответ отсутствуют в представлении сессии и в ответе на проверку. В качестве решения клиент передаёт объект с обязательным строковым полем \texttt{answer}. Пример тела решения:

\begin{verbatim}
{"answer":"  Северный ветер  "}
\end{verbatim}

Проверяющий удаляет пробелы по краям ответа и закреплённого эталона и сравнивает оставшиеся строки. Для эталона \texttt{Северный\ ветер} ответ из примера верен. Строки \texttt{северный\ ветер} и \texttt{Северный\ \ ветер} дадут отрицательный вердикт. Если \texttt{answer} отсутствует, равно \texttt{null} или имеет числовой тип, возвращается \texttt{400 INVALID\_\allowbreak{}REQUEST} без расходования попытки.

\subsection{Результат и отказы}\label{ux440ux435ux437ux443ux43bux44cux442ux430ux442-ux438-ux43eux442ux43aux430ux437ux44b}

\hypertarget{результат-и-отказы}{}

{\def\LTcaptype{none} % do not increment counter
\begin{longtable}[]{@{}
  >{\raggedright\arraybackslash}p{(\linewidth - 2\tabcolsep) * \real{0.3200}}
  >{\raggedright\arraybackslash}p{(\linewidth - 2\tabcolsep) * \real{0.6800}}@{}}
\toprule\noalign{}
\begin{minipage}[b]{\linewidth}\raggedright
Поле результата
\end{minipage} & \begin{minipage}[b]{\linewidth}\raggedright
Тип и правило заполнения
\end{minipage} \\
\midrule\noalign{}
\endhead
\bottomrule\noalign{}
\endlastfoot
\texttt{problemId} & UUID проверенной задачи \\
\texttt{problemType} & Сохранённый тип этой задачи \\
\texttt{isOk} & Логическое значение. \texttt{true} при верном решении, \texttt{false} при неверном \\
\texttt{errorMessage} & Для \texttt{true} --- \texttt{null}. Для \texttt{false} --- непустая строка с пояснением \\
\end{longtable}
}

Оба предметных исхода возвращаются с HTTP 200. При отсутствии профильного пояснения к неверному решению сервер возвращает «Решение неверно». Отказ доступа, ошибка формы, просрочка и технический сбой возвращаются в формате ошибки из следующего раздела. При таком отказе сервер не выдаёт предметный вердикт.

\subsection{Коды и представление ошибок}\label{ux43aux43eux434ux44b-ux438-ux43fux440ux435ux434ux441ux442ux430ux432ux43bux435ux43dux438ux435-ux43eux448ux438ux431ux43eux43a}

\hypertarget{коды-и-представление-ошибок}{}

{\def\LTcaptype{none} % do not increment counter
\begin{longtable}[]{@{}
  >{\raggedright\arraybackslash}p{(\linewidth - 4\tabcolsep) * \real{0.1000}}
  >{\raggedright\arraybackslash}p{(\linewidth - 4\tabcolsep) * \real{0.3400}}
  >{\raggedright\arraybackslash}p{(\linewidth - 4\tabcolsep) * \real{0.5600}}@{}}
\toprule\noalign{}
\begin{minipage}[b]{\linewidth}\raggedright
HTTP
\end{minipage} & \begin{minipage}[b]{\linewidth}\raggedright
Значение \texttt{code}
\end{minipage} & \begin{minipage}[b]{\linewidth}\raggedright
Условие
\end{minipage} \\
\midrule\noalign{}
\endhead
\bottomrule\noalign{}
\endlastfoot
400 & \texttt{INVALID\_\allowbreak{}REQUEST} & Нарушены правила параметров, JSON или обязательных полей \\
401 & \texttt{UNAUTHENTICATED} & Сервер не смог подтвердить личность \\
403 & \texttt{ACCESS\_\allowbreak{}DENIED} & У подтверждённого пользователя отсутствует роль \texttt{USER} \\
404 & \texttt{TASK\_\allowbreak{}NOT\_\allowbreak{}FOUND} & Задачи с таким UUID нет среди собственных \\
404 & \texttt{SESSION\_\allowbreak{}NOT\_\allowbreak{}FOUND} & Сессии с таким UUID нет среди собственных \\
404 & \texttt{ASSESSMENT\_\allowbreak{}NOT\_\allowbreak{}FOUND} & При старте контроля мероприятие отсутствует, ещё не опубликовано или назначено другой группе \\
409 & \texttt{ASSESSMENT\_\allowbreak{}NOT\_\allowbreak{}OPEN} & Старт в своей группе запрошен до \texttt{opensAt} \\
409 & \texttt{ASSESSMENT\_\allowbreak{}CLOSED} & Сервер получил запрос на старт в момент \texttt{closesAt} или позднее \\
409 & \texttt{ASSESSMENT\_\allowbreak{}RETAKE\_\allowbreak{}DISABLED} & Окно открыто, у студента нет сессии \texttt{NEW}, но он уже завершил прохождение этого мероприятия и \texttt{repeatAllowed\ =\ false} \\
404 & \texttt{CONFIGURATION\_\allowbreak{}NOT\_\allowbreak{}FOUND} & У пользователя нет группы или у группы нет включённого шаблона с указанным именем \\
409 & \texttt{ACTIVE\_\allowbreak{}SESSION\_\allowbreak{}EXISTS} & Пользователь уже имеет незавершённый набор \\
409 & \texttt{TASK\_\allowbreak{}FINALIZED} & Своевременный запрос адресован задаче или сессии, которая уже завершилась \\
409 & \texttt{CHECK\_\allowbreak{}IN\_\allowbreak{}PROGRESS} & По этой задаче ещё обрабатывается ранее принятый допустимый ответ \\
410 & \texttt{SESSION\_\allowbreak{}EXPIRED} & Новый ответ к собственной задаче получен ровно в общий срок или позже \\
413 & \texttt{PAYLOAD\_\allowbreak{}TOO\_\allowbreak{}LARGE} & Размер тела больше 1 МиБ \\
500 & \texttt{INTERNAL\_\allowbreak{}ERROR} & Ошибка настройки либо сохранения \\
503 & \texttt{CHECK\_\allowbreak{}UNAVAILABLE} & Проверяющий недоступен или произошёл технический сбой во время его работы \\
503 & \texttt{CHECK\_\allowbreak{}TIMEOUT} & К моменту \texttt{t\_\allowbreak{}in + 60 секунд} результат не зафиксирован \\
\end{longtable}
}

Сначала сервер проверяет доступ. По чужой задаче он возвращает тот же ответ, что и по отсутствующей. Для собственной задачи \texttt{SESSION\_\allowbreak{}EXPIRED} имеет приоритет перед отказом по состоянию. Порядок запуска планировщика не влияет на выбор этой ошибки.

В теле ошибки четыре обязательных строковых поля. Поле \texttt{code} предназначено для выбора действия клиентом, \texttt{message} поясняет причину человеку, \texttt{timestamp} содержит время в том же UTC-формате, что и \texttt{createdAt}, а \texttt{requestId} --- UUID серверной операции. Стек вызовов, текст исключения и служебные данные хранилища клиенту не выдаются. У ошибки \texttt{ACTIVE\_\allowbreak{}SESSION\_\allowbreak{}EXISTS} есть дополнительное поле \texttt{details.\allowbreak{}sessionId} с UUID собственной активной сессии. В ошибке \texttt{ASSESSMENT\_\allowbreak{}RETAKE\_\allowbreak{}DISABLED} поле \texttt{details.\allowbreak{}sessionId} содержит UUID прежней собственной сессии этого мероприятия.

\subsection{Как клиент проходит операции}\label{ux43aux430ux43a-ux43aux43bux438ux435ux43dux442-ux43fux440ux43eux445ux43eux434ux438ux442-ux43eux43fux435ux440ux430ux446ux438ux438}

\hypertarget{как-клиент-проходит-операции}{}

\begin{enumerate}
\def\labelenumi{\arabic{enumi}.}
\tightlist
\item
  При успешном старте клиент сохраняет \texttt{id} сессии и показывает выданные условия. Если ответ на старт контрольной потерян, клиент читает своё расписание через \texttt{GET /\allowbreak{}api/\allowbreak{}users/\allowbreak{}me/\allowbreak{}schedule}. Из карточки этой работы он берёт \texttt{sessionId} своей активной или последней завершённой сессии и открывает её через \texttt{GET /\allowbreak{}api/\allowbreak{}tasks/\allowbreak{}session/\allowbreak{}\{sessionId\}}. Запрос на новый старт используется, когда студент намерен начать прохождение. При самостоятельной практике повторный запрос при наличии действующей сессии даёт \texttt{ACTIVE\_\allowbreak{}SESSION\_\allowbreak{}EXISTS} с её идентификатором в \texttt{details.\allowbreak{}sessionId}.
\item
  Для выбранной задачи клиент отправляет \texttt{PUT} с её UUID и телом по соответствующей схеме. При \texttt{400 INVALID\_\allowbreak{}REQUEST} можно исправить тело и отправить его до общего срока. Предметный ответ с \texttt{isOk\ =\ false} завершает попытку.
\item
  Если после отправки соединение прервалось, клиент читает сессию по её UUID, чтобы узнать, что зафиксировал сервер.
\end{enumerate}

\subsection{Принятая обработка и восстановление}\label{ux43fux440ux438ux43dux44fux442ux430ux44f-ux43eux431ux440ux430ux431ux43eux442ux43aux430-ux438-ux432ux43eux441ux441ux442ux430ux43dux43eux432ux43bux435ux43dux438ux435}

\hypertarget{принятая-обработка-и-восстановление}{}

«Принятый ответ» --- запрос с решением владельца, который сервер уже зарегистрировал для обработки. Получение запроса ещё не означает, что решение верно или что вердикт сохранён. Сервер различает {\hyperlink{время-решения-и-время-обработки}{время подачи и предел обработки}}. Например, при подаче за одну секунду до \texttt{closesAt} предельное время записи вердикта наступит через 59 секунд после закрытия. Новый ответ владельца в момент \texttt{deadlineAt} получит \texttt{410 SESSION\_\allowbreak{}EXPIRED}.

При полном получении запроса сервер отмечает \texttt{t\_in}. До очереди он сверяет подтверждённого отправителя с владельцем задачи. Чужой или отсутствующий UUID даёт \texttt{404 TASK\_\allowbreak{}NOT\_\allowbreak{}FOUND} и не открывает операцию для этой задачи. Для собственной задачи сервер записывает в устойчивое хранилище UUID операции, владельца, UUID задачи, полный ответ, \texttt{t\_in}, предел \texttt{t\_\allowbreak{}in + 60 секунд} и признак незавершённости.

Своевременно зарегистрированный ответ владельца защищает задачу от просрочки во время ожидания, допуска и проверки формата. При отказе защита по этому запросу снимается. Право начать предметную проверку подсистема предоставляет атомарно для задачи. Если два допустимых запроса зарегистрированы до срока, при идущей проверке второй получает \texttt{409 CHECK\_\allowbreak{}IN\_\allowbreak{}PROGRESS}, после фиксации итога до общего срока --- \texttt{409 TASK\_\allowbreak{}FINALIZED}. Повторный запрос и перезапуск приложения не заменяют исходные данные и не отодвигают предел.

После перезапуска сервер должен продолжить обработку открытой записи по сохранённому ответу, пока исходный предел не достигнут. Если вердикт уже зафиксирован, сервер выдаёт его при чтении. Если предел наступил до фиксации, операция завершается с \texttt{CHECK\_\allowbreak{}TIMEOUT}. Пока общий срок не истёк, пользователь может отправить новый ответ. После общего срока задача без своевременной обработки закрывается по просрочке.

Если при записи вердикта произошёл откат, сервер сообщает клиенту \texttt{500 INTERNAL\_\allowbreak{}ERROR} только после сохранения отмены операции. Прежний вердикт после такого отказа уже не появится. Если вместо ответа клиент обнаружил разрыв соединения, он читает собственную сессию. При открытой записи сервер продолжит обработку в пределах исходного времени, при уже сохранённом вердикте вернёт результат чтением.

\subsection{Чтение после потери ответа}\label{ux447ux442ux435ux43dux438ux435-ux43fux43eux441ux43bux435-ux43fux43eux442ux435ux440ux438-ux43eux442ux432ux435ux442ux430}

\hypertarget{чтение-после-потери-ответа}{}

Потеряться может HTTP-ответ сервера с вердиктом даже тогда, когда сам вердикт уже записан. Для восстановления клиент читает личную сессию по её UUID. Обозначение \texttt{D} ниже означает общий срок подачи решений.

Операция чтения доступна владельцу до и после \texttt{D}. Перед выдачей сервер учитывает просрочку и сохраняет право на завершение своевременного ответа в обработке. Клиент различает исходы по сочетанию \texttt{state} и \texttt{result}.

{\def\LTcaptype{none} % do not increment counter
\begin{longtable}[]{@{}
  >{\raggedright\arraybackslash}p{(\linewidth - 2\tabcolsep) * \real{0.3200}}
  >{\raggedright\arraybackslash}p{(\linewidth - 2\tabcolsep) * \real{0.6800}}@{}}
\toprule\noalign{}
\begin{minipage}[b]{\linewidth}\raggedright
Задача в ответе на чтение
\end{minipage} & \begin{minipage}[b]{\linewidth}\raggedright
Что произошло
\end{minipage} \\
\midrule\noalign{}
\endhead
\bottomrule\noalign{}
\endlastfoot
\texttt{NEW}, \texttt{result\ =\ null} & Предметного вердикта пока нет. Своевременный запрос может ещё обрабатываться \\
\texttt{SOLVED}, в результате \texttt{isOk\ =\ true} & Зафиксирован верный ответ \\
\texttt{FAILED}, в результате \texttt{isOk\ =\ false} & Зафиксирован неверный ответ. Причина приведена в \texttt{errorMessage} \\
\texttt{FAILED}, \texttt{result\ =\ null} & Задача завершена по сроку без предметного вердикта \\
\end{longtable}
}

Чтение выдаёт сохранённые результаты без повторного вызова проверяющего и без расхода попытки. Пользователь, запросивший чужую сессию, получает \texttt{404 SESSION\_\allowbreak{}NOT\_\allowbreak{}FOUND}, как и при отсутствии указанного UUID.


\end{document}
